% SOURCE_AUTHORITY: sga5_fr_workpass.tex, lines 4441--4494 (LF-normalized unit).
% SOURCE_SHA256: 791F4EFFC5E02832D5D77ED03518C8156D6F07E4C8238B03545DB93D883FBB28
La idea de la demostración es que, en este caso, $u$ y $v$ son imágenes directas de correspondencias entre complejos concentrados en los puntos cerrados y que la fórmula que debe establecerse no hace sino expresar la fórmula de Lefschetz para las proyecciones $x\to S$ e $y\to S$. Teniendo en cuenta 1.1.3, se pueden sustituir los datos de 1.1 por los siguientes: $X$ (resp. $Y$) es una $S$-curva lisa provista de un punto cerrado $x$ (resp. $y$), $A$ (resp. $B$) es una $S$-curva separada, y $a:A\to X\times_S Y$, $b:B\to X\times_S Y$ son unos $S$-morfismos finitos tales que $A\times_{X\times_SY}B$ se reduce a un punto cerrado $z$ situado sobre $(x,y)$,
\[
u\in\Hom(a_1^*L,a_2^!M),\qquad v\in\Hom(b_2^*M,b_1^!L),
\]
con
\[
L=h_{1*}L',\quad M=h_{2*}M',\quad
L'\in\Ob D_{\mathrm{ctf}}(\{x\}),\quad M'\in\Ob D_{\mathrm{ctf}}(\{y\}),
\]
donde $h_1:\{x\}\to X$ y $h_2:\{y\}\to Y$ designan las inclusiones canónicas. Se trata de demostrar que
\[
\langle u_z,v_z\rangle=\langle u,v\rangle_z,
\]
donde
\[
u_z\in\Hom(\mathrm R\Gamma(X,L),\mathrm R\Gamma(Y,M)),\qquad
v_z\in\Hom(\mathrm R\Gamma(Y,M),\mathrm R\Gamma(X,L))
\]
están definidos como en 1.1, y $\langle u,v\rangle_z$ es el valor en $z$ del producto cup (III 4.2.7). Se tiene $\mathrm R\Gamma(X,L)=\mathrm R\Gamma_c(X,L)=L_x$, $\mathrm R\Gamma(Y,M)=\mathrm R\Gamma_c(Y,M)=M_y$, y $u_z$ (resp. $v_z$) no es sino la imagen directa $u_*$ (resp. $v_*$) de $u$ (resp. $v$) sobre $S$ en el sentido de (III 3.7.6), como muestra el procedimiento de cálculo (III 3.5 b)).

Por otra parte, si $h=h_1\times_Sh_2:(x,y)\to X\times_SY$ es la inclusión, y si $a':A'\to(x,y)$ y $b':B'\to(x,y)$ designan las fibras de $a$ y $b$ en $(x,y)$, cada una de ellas reducida a un punto, se tienen isomorfismos canónicos
\[
h_*a'_*a'{}^!(DL_x\lten_SM_y)
 \xrightarrow{\sim} a_*a^!(DL\lten_SM),
\]
\[
h_*b'_*b'{}^!(L_x\lten_SDM_y)
 \xrightarrow{\sim} b_*b^!(L\lten_SDM),
\]
y las flechas
\[
h_*:\Hom(a_1'{}^*L_x,a_2'{}^!M_y)\to \Hom(a_1^*L,a_2^!M),
\]
\[
h_*:\Hom(b_2'{}^*M_y,b_1'{}^!L_x)\to \Hom(b_2^*M,b_1^!L)
\]
(III 3.7.6) son isomorfismos. Por tanto, las correspondencias $u$ y $v$ se escriben de manera única
\begin{equation}\tag{1}
u=h_*u',\qquad v=h_*v',
\end{equation}
con
\[
u'\in\Hom(a_1'{}^*L_x,a_2'{}^!M_y),\qquad
v'\in\Hom(b_2'{}^*M_y,b_1'{}^!L_x).
\]
Por la fórmula de Lefschetz (III 4.5.1) para $(h_1,h_2,A'\to A,B'\to B)$, se tiene
\begin{equation}\tag{2}
\langle u,v\rangle_z=\langle u',v'\rangle_z .
\end{equation}
Por otra parte, según (1), se tiene $u_*=u'_*$, $v_*=v'_*$, de donde, por la fórmula de Lefschetz para $\{x\}\to S$, $\{y\}\to S$,
\begin{equation}\tag{3}
\langle u_*,v_*\rangle=\langle u',v'\rangle_z .
\end{equation}
Teniendo en cuenta la observación anterior, la conclusión resulta al combinar (2) y (3).
